\subsection{The Nyman--Beurling comparison}
\label{subsec:mobius-nyman}

The Mellin inverse also gives a precise comparison with the
Nyman--Beurling approximation problem. Put
\[
 \chi(x)=\mathbf1_{(0,1)}(x),\qquad
 r_n(x)=\{1/(nx)\},\qquad
 D_N(s)=\sum_{n\le N}a_n n^{-s}.
\]
The coefficient map \(a\mapsto\chi+\sum_{n\le N}a_n r_n\)
satisfies
\begin{equation}\label{eq:nyman-mellin-map}
 \left\|\chi+\sum_{n\le N}a_n r_n\right\|_{L^2(0,\infty)}^2
 =\frac1{2\pi}\int_{\mathbb R}
  \frac{|1-\zeta(1/2+it)D_N(1/2+it)|^2}{1/4+t^2}\,dt.
\end{equation}
Indeed, for \(0<\Re s<1\), integration of the integer counting
function, followed by \(t=1/x\), gives
\[
 \int_0^\infty r_1(x)x^{s-1}\,dx=-\frac{\zeta(s)}s.
\]
One can see the boundary terms explicitly: the part \(0<t<1\)
of \(\int_0^\infty\{t\}t^{-s-1}\,dt\) is \(1/(1-s)\),
and the part \(t\ge1\) is \(1/(s-1)-\zeta(s)/s\).
They add to the displayed formula. Dilation gives the factor
\(n^{-s}\), whereas the Mellin transform of \(\chi\) is
\(1/s\). Mellin Plancherel proves \eqref{eq:nyman-mellin-map}.

Thus the squared distance \(d_N^2\), defined as the infimum in
\eqref{eq:nyman-mellin-map}, is a finite Gram minimum. Its entries
are
\[
 G_{nm}=\int_0^\infty r_n(x)r_m(x)\,dx,\qquad
 b_n=\int_0^1r_n(x)\,dx
       =\frac{1-\gamma_{\!E}+\log n}{n},
 \qquad d_N^2=1-b^TG^{-1}b.
\]
Here \(\gamma_{\!E}\) is Euler's constant. To evaluate \(b_n\),
substitute \(y=nx\); the integral on \((1,n)\) is \(\log n\),
and integration over the intervals
\((1/(j+1),1/j)\) gives
\(\int_0^1\{1/y\}\,dy=1-\gamma_{\!E}\).
The matrix \(G\) is positive definite. In a vanishing linear
combination of the \(r_n\), the jump at \(x=1\) first gives
\(a_1=0\). At \(x=1/j\), the jumping terms have indices
dividing \(j\); induction gives \(a_j=0\) for \(1\le j\le N\).
A function vanishing almost everywhere has zero jump because it
vanishes on both adjacent open intervals. Completing the square
now proves the formula for \(d_N^2\).
C\'aceres \cite[Sections~14, 16--17, pp.~50--52]{caceres-margin-v2}
also compares this finite Gram minimum with the zero-sum constant.
His linear coefficient is exactly \(c_n=-b_n\), and his Gram
matrix is the same fractional-part autocorrelation matrix; hence
\(1-c^TG^{-1}c=1-b^TG^{-1}b\).


Burnol's lower bound \cite[Theorem~1.3, p.~4; Section~5,
pp.~14--16]{burnol2002} yields
\begin{equation}\label{eq:nyman-multiplicity-excess}
 \liminf_{N\to\infty}(\log N)d_N^2
 \ge 2+\gamma_{\!E}-\log(4\pi)
   +\sum_{\rho\ {\rm distinct}}
       \frac{m_\rho(m_\rho-1)}{|\rho|^2}.
\end{equation}
The sum includes both signs of the ordinate. To specify the
comparison, each generator \(r_n\) in \eqref{eq:nyman-mellin-map}
is Burnol's fractional-part generator at parameter \(\theta=1/n\).
His continuous-parameter approximation space on \([1/N,1]\)
therefore contains this finite space, so its distance is at most
\(d_N\). His theorem gives the lower sum
\(\sum m_\rho^2/|\rho|^2\). The paired Hadamard product and the
functional equation give
\[
 \sum_\rho\frac{m_\rho}{|\rho|^2}
 =\frac{\xi'}\xi(1)-\frac{\xi'}\xi(0)
 =2+\gamma_{\!E}-\log(4\pi),
\]
where the zero-location theorem identifies
\(\rho(1-\rho)=|\rho|^2\). Subtracting this identity proves
\eqref{eq:nyman-multiplicity-excess}. The sums converge: the
local zero count and multiplicity bound give an upper bound
\(O((\log(2+|\gamma|))^2)\) for the squared multiplicities in
each unit interval, which is summable after division by
\(1/4+\gamma^2\).

The squared multiplicity in this formula comes from an explicit
jet Gram matrix. Burnol's normalized vectors at a zero of
multiplicity \(m\) have limiting matrix
\((1/(j+k+1))_{0\le j,k<m}\); its inverse has upper left entry
\(m^2\) \cite[Theorems~5.3--5.5, pp.~14--16]{burnol2002}.
This is the norm of the evaluation functional at zero on the
local jet space after the indicated rescaling.

The dependence on those jets can already be computed for one
double zero. In the Hardy space on the unit circle, let
\(P(z)=(1-z)^2\) and taper the coefficients of
\(P^{-1}=\sum_{n\ge0}(n+1)z^n\):
\[
 W_N(z)=\sum_{n=0}^N(1-n/N)(n+1)z^n,\qquad N\ge1.
\]
Multiplication gives the exact residual
\[
 1-P(z)W_N(z)=\frac2N\sum_{n=1}^Nz^n-z^{N+1},
 \qquad \|1-PW_N\|_{H^2}^2=1+\frac4N.
\]
The optimal degree-\(N\) multiplier has the smaller exact error
\begin{equation}\label{eq:nyman-double-example}
 \min_{\deg Q\le N}\|1-PQ\|_{H^2}^2
       =\frac{2(2N+5)}{(N+3)(N+4)}.
\end{equation}
To verify it, write the residual as a polynomial of degree at
most \(M=N+2\). Divisibility by \((1-z)^2\) says precisely
that the residual has value one and derivative zero at \(1\).
The two representing vectors have Gram matrix
\[
 \begin{pmatrix}
 M+1&M(M+1)/2\\
 M(M+1)/2&M(M+1)(2M+1)/6
 \end{pmatrix}.
\]
The squared minimum is the upper left entry of its inverse,
\(2(2M+1)/((M+1)(M+2))\), giving
\eqref{eq:nyman-double-example}. This is the double-zero case
of the polynomial approximation mechanism in
\cite[Section~2, pp.~4--6]{burnol2002}.

For zeta, Bettin--Conrey--Farmer use the logarithmically tapered
M\"obius polynomial
\(V_N(s)=\sum_{n\le N}(1-\log n/\log N)\mu_{\mathrm M}(n)n^{-s}\)
\cite[Theorem~1, p.~2; Section~3, pp.~5--9]{bettin-conrey-farmer2013}.
Their upper asymptotic uses a bound on the full sum of
\(|\zeta'(\rho)|^{-2}\). A multiple zero makes its summand
infinite, so finiteness of that input already includes simplicity
on each bounded packet. The local pole formulas above display
the corresponding higher-order terms before any such bound is
available. An upper estimate matching the first constant in
\eqref{eq:nyman-multiplicity-excess} has not been proved here.
